% ──────────────────────────────────────────────────────────────────────────────
%  Section 1 — Setup
% ──────────────────────────────────────────────────────────────────────────────

\section{Setup and line-only constants}
\label{sec:setup}

The line is modelled in the distributed-parameter form developed in the
parent technical document. The local terminal voltage is held balanced
and positive sequence at the nominal line-to-neutral magnitude
$V_{s} = V_{LL}/\sqrt{3}$, with phase ordering
$\vu = (1,\,a^{2},\,a)^{\top}$ and $a = \ee^{\jj 2\pi/3}$. The
sending-end three-phase complex power and the line length are the two
free parameters of the present analysis,
$S = P + \jj Q \in \mathbb{C}$ and
$\ell \in \mathbb{R}_{+}$.

The line geometry enters the equations only through the matrix
exponential of the state-space generator,
$\matM(\ell) = \ee^{\matF \ell}$, with
$\matF = [\,\bm{0},\,-\matZ;\,-\matY,\,\bm{0}\,]$ assembled from the
per-unit-length series impedance and shunt admittance matrices. The
upper blocks of $\matM(\ell)$ are the ABCD operators $\matA(\ell)$ and
$\matB(\ell)$, both of which are independent of $S$ and $V_{s}$. From
these blocks four \emph{line-only} quantities are formed, all of which
will be repeatedly invoked below:
%
\begin{equation}
\label{eq:line-constants}
\vp(\ell) \;=\; \matB(\ell)^{-1}\,\vu,
\qquad
\vq(\ell) \;=\; \matB(\ell)^{-1}\matA(\ell)\,\vu,
\end{equation}
%
\begin{equation}
\label{eq:alpha-beta}
\alpha(\ell) \;=\; \vu \cdot \overline{\vp(\ell)},
\qquad
\beta(\ell)  \;=\; \vu \cdot \overline{\vq(\ell)},
\end{equation}
%
\noindent where the dot product is the ordinary, non-Hermitian component
sum and the overline denotes complex conjugation. The vectors
$\vp,\vq \in \mathbb{C}^{3}$ carry the phase index $k \in \{A,B,C\}$
through their components $p_{k}$ and $q_{k}$; the scalars
$\alpha,\beta \in \mathbb{C}$ do not. Every quantity in
\eqref{eq:line-constants}--\eqref{eq:alpha-beta} depends on $\ell$ only
and is computed once for a given tower geometry by exponentiating
$\matF$ at the requested length.

The two-point boundary problem solved in the parent document delivers
the receiving-end complex unknown
\begin{equation}
\label{eq:W-recall}
W(\ell, S)
\;=\;
\frac{S \;+\; V_{s}^{2}\,\beta(\ell)}{V_{s}\,\alpha(\ell)},
\end{equation}
\noindent and the sending-end phase currents follow from
$\vI_{s} = \matB^{-1}\!\bigl(W^{*}\vu - V_{s}\,\matA\vu\bigr)$. Both
relations are recalled here only to fix notation; the rest of the
present note works directly in $(S,\ell)$.
